% !TeX root = ../main.tex

\chapter{Analisi Economica: Costi, ROI e ROA (Return on Agent)}
\label{cap:analisi_economica}

\section{Service Level Objectives (SLO) e Costo del Downtime}
\label{sec:slo_downtime}

\subsection*{Definizione degli SLO}
Per analizzare l'affidabilità del sistema e definirne la tolleranza ai guasti, sono stati stabiliti tre \textit{Service Level Objectives} (SLO), coerenti con il workload presentato nei capitoli precedenti:

\begin{itemize}
    \item \textbf{SLO 1 -- Latenza di Assegnazione Logistica:} il 90\% degli ordini ad alta priorità deve essere processato e assegnato a un mezzo entro 1 minuto. Il valore è giustificato dal ciclo di vita dell'orchestratore, che prevede un \texttt{time.sleep(30)} di base a cui si somma la latenza di inferenza dell'API.
    \item \textbf{SLO 2 -- Disponibilità del Data Layer:} il broker MQTT e il database InfluxDB devono garantire il 99{,}9\% di uptime. Considerato che i droni inviano telemetria ogni 2 secondi, una perdita prolungata di dati (per quanto mitigata dalle tecniche viste nel Capitolo~\ref{cap:reliability_scalability}) renderebbe cieca l'IA, paralizzando l'\texttt{HealthAgent} nella sua funzione di monitoraggio con possibili risvolti pericolosi.
    \item \textbf{SLO 3 -- Affidabilità e risoluzione agentica:} l'LLM deve concludere il proprio task di ragionamento ed emettere una \textit{tool call} valida nel 99\% dei casi, senza innescare l'interruttore di sicurezza \texttt{MAX\_AGENT\_ITERATIONS}, evitando timeout e spreco di computazione.
\end{itemize}

\subsection*{Stima del Costo di un'Ora di Downtime}
Operando nel dominio della logistica critica, il costo stimato di un'ora di downtime del servizio è di circa \textbf{15.000\,\texteuro/ora}. La stima deriva da una media dei guadagni che il servizio raggiunge in un'ora di attività portando a termine le consegne, integrata dalle possibili penali contrattuali per violazione degli SLA.

\section{Valutazione del Return on Agent (ROA) e Costi}
\label{sec:roa}

L'orchestrazione autonoma della flotta introduce un bilanciamento critico tra i costi operativi delle API (OPEX) e il risparmio di capitale hardware (CAPEX) ottenuto tramite logiche predittive.

\subsection{Strategie di Ottimizzazione dei Costi Operativi (OPEX)}
Il sistema limita l'impatto economico delle chiamate LLM attraverso tre strategie integrate:

\begin{itemize}
    \item \textbf{Triage preventivo:} il modulo \texttt{triage\_manager} analizza lo snapshot del sistema prima di invocare i modelli, attivando \texttt{LogisticAgent} ed \texttt{HealthAgent} solo se sussistono reali condizioni di sbilanciamento (es. ordini pendenti accoppiati a droni idle, o necessità di scaling/manutenzione). Questo approccio azzera i costi derivanti da cicli di no-op ridondanti.
    \item \textbf{Ingegneria dei prompt:} i prompt di sistema impongono restrizioni sintattiche che vietano la generazione di testo libero, preamboli o spiegazioni discorsive del ragionamento. L'output dei modelli è confinato alle sole chiamate di funzione (\textit{tool calling}), minimizzando drasticamente il consumo di token di completamento.
    \item \textbf{Monitoraggio dei consumi:} il sistema esegue un auditing sistematico dei token spesi (prompt, completamento e totale) al termine di ogni esecuzione, permettendo una precisa rendicontazione finanziaria del processo decisionale.
\end{itemize}

\subsection{Mitigazione del Rischio e Preservazione del Capitale (CAPEX)}
Il ritorno sull'investimento dell'agente intelligente si concretizza nella mitigazione dei rischi d'impresa e nell'ottimizzazione infrastrutturale:

\begin{itemize}
    \item \textbf{Manutenzione preventiva dinamica:} il \texttt{LogisticAgent} include nei prompt operativi la regola di evitare l'assegnazione di carichi pesanti ($>3$\,kg) a droni con usura superiore al $70\%$. Questa policy, pur essendo oggi principalmente prompt-driven (\textit{soft constraint}), orienta il matching ordine-drone verso scelte conservative che riducono il rischio di esaurimento batteria e perdita del vettore.
    \item \textbf{Elasticità infrastrutturale cloud:} l'\texttt{HealthAgent} modula dinamicamente le repliche del Deployment \texttt{drone-simulator} su Kubernetes basandosi sulle fluttuazioni del carico di lavoro effettivo, riducendo i costi di over-provisioning dell'infrastruttura cloud.
    \item \textbf{Operational Safety Cap:} le policy di sicurezza impongono un limite massimo di 6 droni per lo scaling automatico. Qualsiasi espansione ulteriore viene intercettata e subordinata ad approvazione umana (\texttt{request\_human\_approval}), ponendo un vincolo rigido (\textit{budget cap}) contro anomalie o loop di scaling incontrollati.
    \item \textbf{Massimizzazione del throughput:} la coda degli ordini viene ordinata preventivamente per priorità. L'agente logistico alloca la capacità di carico residua privilegiando i task \textit{high priority}, lasciando il valore economico come criterio secondario nel ragionamento del modello.
\end{itemize}

\label{sec:opex_budget}

\subsection*{Proiezioni Economiche e Budget Compliance}
Per convalidare la sostenibilità finanziaria del software è stata effettuata una stima comparativa dei costi su base mensile, ipotizzando una frequenza di esecuzione standard di 2.880 cicli giornalieri (un polling di verifica ogni 30 secondi).

\begin{table}[ht]
\centering
\begin{tabular}{lrr}
\toprule
\textbf{Voce di Costo} & \textbf{Gemini 2.5 Pro} & \textbf{Mistral Small} \\
\midrule
Infrastruttura Cloud (1 control-plane + 2 worker) & 150{,}00 \$ / mese & 150{,}00 \$ / mese \\
Token Modello Linguistico (AI) & $\sim$500{,}00 \$ / mese & $\sim$10{,}00 \$ / mese \\
\midrule
\textbf{OPEX Totale Mensile} & \textbf{$\sim$650{,}00 \$ / mese} & \textbf{$\sim$160{,}00 \$ / mese} \\
\bottomrule
\end{tabular}
\caption{Confronto dei costi mensili al variare del modello LLM utilizzato.}
\label{tab:costi_opex}
\end{table}

\begin{itemize}
    \item \textbf{Costi infrastrutturali fissi:} il mantenimento dell'ambiente containerizzato in alta affidabilità (comprensivo di 1 control-plane e 2 nodi worker Kubernetes, database InfluxDB per la telemetria e broker MQTT per la messaggistica) genera un costo operativo fisso di circa \textbf{150\,\$/mese}.
    \item \textbf{Costi variabili dell'Intelligenza Artificiale:} l'impatto della componente generativa è fortemente condizionato dall'architettura del modello scelto. \textit{Gemini 2.5 Pro} comporta un volume medio di 1.200 token per ciclo (ripartiti equamente tra input e output), traducendosi in una spesa di circa \textbf{500\,\$/mese}. Un modello specialistico più leggero come \textit{Mistral Small} riduce l'impronta a 800 token per ciclo; grazie a un'architettura asimmetrica focalizzata sull'input (90\% input, 10\% output), il costo computazionale si abbatte a meno di \textbf{10\,\$/mese}.
\end{itemize}

L'analisi finanziaria evidenzia che la quasi totalità del costo variabile è legata al modello linguistico sottostante. Mentre l'infrastruttura cloud mantiene un costo fisso e invariabile di 150\,\$/mese, la spesa per l'intelligenza artificiale fluttua significativamente, generando una marcata \textbf{forbice dei costi} che spazia dai 500\,\$/mese di \textit{Gemini 2.5 Pro} a meno di 10\,\$/mese per \textit{Mistral Small}. Questo divario si traduce in un delta operativo di circa \textbf{490\,\$ mensili}, offrendo la possibilità di modulare la scelta dell'LLM in base al livello di complessità decisionale richiesto in tempo reale.

Per scenari ordinari di instradamento degli ordini e scaling lineare, la flotta può essere orchestrata efficacemente da modelli specialistici e leggeri, riducendo l'OPEX totale ad appena \textbf{160\,\$/mese}. Al contrario, l'allocazione del modello top-tier può essere riservata esclusivamente a finestre temporali critiche caratterizzate da anomalie infrastrutturali complesse, elevata congestione della coda o conflitti di policy che richiedono capacità di ragionamento superiori.

\section{Analisi dei Trade-Off}
\label{sec:analisi_tradeoff}

\subsubsection*{1. Trade-off tra Costo e Affidabilità}
L'adozione di un sistema basato su agenti LLM introduce un'inversione di paradigma rispetto a un sistema deterministico tradizionale, ridefinendo il concetto stesso di continuità operativa:

\begin{itemize}
    \item \textbf{Sistema deterministico:} presenta costi di computazione (OPEX) minimi, limitati all'esecuzione di script locali. Il limite principale risiede nella resilienza operativa: a fronte di eccezioni non gestite o anomalie di rete, i flussi tradizionali subiscono arresti critici. Il ripristino del servizio richiede l'intervento manuale di un amministratore di sistema, incrementando i costi legati alla reperibilità oraria e introducendo latenze nei tempi di risoluzione.
    \item \textbf{Sistema multi-agente:} introduce un costo operativo variabile legato al consumo di token delle API (stimato fino a 500\,\$/mese con \textit{Gemini 2.5 Pro}). Questo investimento si traduce in una migliore gestione degli imprevisti grazie alla flessibilità semantica del modello, che gli consente di contestualizzare l'errore, ragionare sulla deviazione e deviare dinamicamente su flussi alternativi, garantendo la continuità del servizio senza richiedere l'intervento umano.
\end{itemize}

\subsubsection*{2. Trade-off tra Automazione e Sicurezza (Automation vs. Safety)}
Questo bilanciamento definisce le soglie di autonomia decisionale del software prima della transizione del controllo all'operatore umano:

\begin{itemize}
    \item \textbf{Sistema deterministico:} garantisce l'osservanza rigorosa delle policy di sicurezza e dei vincoli operativi grazie a logiche cablate nel codice. Tale rigidità ne limita tuttavia la flessibilità: a fronte di anomalie o picchi improvvisi di carico (come la sovrapposizione tra l'accumulo di ordini prioritari e il fermo contemporaneo di più vettori) l'algoritmo non è in grado di rimodulare autonomamente le soglie di tolleranza o di deviare dai flussi standard, provocando il blocco dei processi logistici fino a un intervento manuale.
    \item \textbf{Sistema multi-agente:} massimizza l'efficienza dei processi decisionali, abilitando gli agenti a valutare trade-off complessi tra metriche di business (ottimizzazione dei ricavi) e vincoli fisici (usura dei vettori). Di contro, l'elasticità decisionale introduce rischi legati alla sicurezza deterministica, quali potenziali allucinazioni nell'invocazione dei tool o derive logiche nelle richieste di scaling infrastrutturale, mitigati dai guardrail descritti nel Capitolo~\ref{cap:safety_guardrails}.
\end{itemize}
